% ============================================
% DOC-03 (ZH)
% 中文翻译 — SPEC-935-R356
% ============================================
\section*{\textbf{DOC-03} --- 尸检报告}
\markboth{DOC-03 --- 尸检报告}{}

\noindent \textit{巴尔巴塞纳法医研究所 —— 1979年6月14日 —— 阿尔瓦罗·托雷斯医生}

\noindent \textbf{案卷编号:} IML-BA 1979.06.12-1260
\textbf{死者:} 若阿金·安东尼奥·科雷亚（医院登记号1,260）
\textbf{死亡日期:} 1979年6月12日
\textbf{法医:} 阿尔瓦罗·托雷斯医生，CRM-MG 3,117

\vspace{0.5em}\noindent\rule{\textwidth}{0.5pt}\vspace{0.5em}

\noindent \textbf{巴尔巴塞纳法医研究所 —— 尸检报告第1,260号 —— 1979年6月14日}

\noindent \textbf{死者:} 若阿金·安东尼奥·科雷亚
\textbf{登记号:} 1,260
\textbf{估计年龄:} 72岁（±5岁）
\textbf{籍贯:} 塞阿拉州基谢拉莫宾
\textbf{死亡日期:} 1979年6月12日
\textbf{尸检日期:} 1979年6月14日
\textbf{法医:} 阿尔瓦罗·托雷斯医生 —— CRM-MG 3,117
\textbf{委托方:} 巴尔巴塞纳收容医院领导层

\vspace{0.5em}\noindent\rule{\textwidth}{0.5pt}\vspace{0.5em}

\subsection*{1. 外部检查}

\subsubsection*{1.1 人体测量数据}

\noindent - \textbf{体重:} 42公斤（身体质量指数：16.0公斤/平方米 —— 严重慢性营养不良）
- \textbf{身高:} 1.62米
- \textbf{营养状况:} 明显恶病质。皮下脂肪缺失。全身肌肉量减少。骨骼突起——尤其是锁骨、髂嵴和髌骨——突出于皮下，骨性隆起显著。皮肤干燥，皮肤皱褶持续存在（慢性脱水征象）。指甲呈现博氏线（横向沟纹，提示一生中反复发作的急性营养不良）。

\subsubsection*{1.2 眼部发现}

\noindent \textbf{单侧不对称黄疸 —— 右眼球。}

\noindent 右眼巩膜呈强烈的金黄色，异乎寻常，从角膜延伸至球结膜，色调均匀。该颜色既不符合典型肝性黄疸（那会是双侧弥漫性的），也不符合眼睑黄色瘤。左眼巩膜呈白色，血管分布正常。

\noindent \textit{法医注:} 触诊时，右眼球质地比左眼略硬。无近期创伤迹象。无假体。黄色巩膜在指压后不褪色。

\noindent \textit{假设:}
- 未经治疗的眼部创伤后遗症
- 未记录的新生儿感染
- \textit{（页边批注，不同笔迹：“掘尸之后，这一假设不成立了。眼睛在黑暗中发光。这一现象没有医学解释。 —— 卢西亚·M.”）}

\subsubsection*{1.3 电休克治疗（ECT）损伤}

\noindent 颅骨在双侧颞部、双耳屏正前方呈现\textbf{双侧电灼伤痕迹}。痕迹呈椭圆形，约3厘米×2厘米，褐黄色，边缘略隆起，瘢痕组织增厚。这些区域的皮肤呈永久性脱发——毛囊已被反复通过的电流破坏。

\noindent \textit{分析:} 这些痕迹符合在没有导电凝胶的情况下反复使用干性金属电极进行电休克治疗的特征。导电凝胶的缺失——收容院直至1960年代的标准做法——造成部分厚度烧伤，愈合后破坏了该区域的汗腺和毛囊。

\noindent \textbf{与ECT相符的骨折:}
- \textbf{T5椎体压缩性骨折}（第五节胸椎），椎体高度损失约30\%。陈旧性愈合，无过度骨痂。典型于全身性强直阵挛发作期间剧烈、失控的肌肉收缩——手术时患者仰卧，无适当约束，强直期脊柱的强制屈曲将椎体压缩至骨折。
- \textbf{左股骨颈骨折已愈合}，左下肢缩短约2厘米（据护理记录，行走时观察到）。骨折发生于1948年一次电休克治疗中——患者在阵挛期被抛出担架，摔在髋部。
- \textbf{右下颌支骨折已愈合}，面部轻度不对称。抽搐发作期间，咬肌收缩力量足以使下颌骨骨折。缺少口腔保护（护齿或橡胶卷）促成了该损伤。

\subsubsection*{1.4 经眶额叶切除术（弗里曼术式）痕迹}

\noindent \textbf{双眼内眦}均有约2毫米直径的点状瘢痕，位于泪阜与半月皱襞交界处。瘢痕双侧对称。

\noindent \textit{分析:} 瘢痕符合经眶白质刀的导入——一种类似碎冰锥的器械，长约15厘米，直径6毫米，钝头，柄上有刻度。该术式由沃尔特·弗里曼于1946年描述，步骤如下：

\noindent 1. 先对患者施行电休克以诱导意识丧失（“惊厥麻醉”）
2. 将白质刀置于眼球结膜上，位于眼内眦
3. 将器械向中线倾斜20度
4. 用手术锤叩击白质刀，直至穿透眶骨薄板（约2—3毫米厚）
5. 将器械推进至颅腔内5—7厘米
6. 呈扇形移动白质刀柄（由内向外、由外向内），切断前额叶联合纤维
7. 对侧重复操作

\noindent \textit{相关发现:} 在筛板（眶顶）上可见两个约4毫米直径的小孔，边缘完全光滑且已愈合，对称地位于距中线2.5厘米处。孔间距与骨缘角度证实器械是朝向内侧额叶方向导入的。\textbf{孔道通畅} —— 未见完全骨性愈合。

\noindent \textit{法医页边注:} “尽管骨上孔道明显，大脑肉眼检查未发现与手术相符的皮质瘢痕。额叶切除术可能未完成，或者器械沿一条未对脑组织造成可检测损伤的解剖路径行进。注：病历中没有任何额叶切除术授权。该手术未留下正式记录。”

\subsubsection*{1.5 身体约束痕迹}

\noindent \textbf{手腕与脚踝:} 双腕（宽3厘米）与双踝（宽2.5厘米）均有环周色素沉着带。这些区域的皮肤粗糙，失去弹性，无毛发。带痕连续均匀，为\textbf{皮革手铐或约束衣长期压迫}的特征。

\noindent \textit{观察:} 手腕带痕深到足以造成双侧正中神经压迫（约束性腕管综合征）。手指半屈曲挛缩（爪形手）提示患者多次被连续约束超过30天。

\noindent \textbf{胸腹:} 乳头高度有双侧水平线性擦伤，符合皮条或帆布带（约束衣）反复摩擦的特征。擦伤愈合程度不一，提示数十年间反复使用。

\noindent \textbf{髋部与大腿:} 双大腿前侧、腹股沟褶皱高度有脱发与真皮纤维化区域。符合长期卧床制动——患者仰卧，双腿由皮条或湿床单约束（床单干后紧贴皮肤，造成组织缺血）。

\subsubsection*{1.6 自残瘢痕}

\noindent 双前臂前后面有多条线性瘢痕。分布对称，时间不一（从不到6个月的新瘢痕到数十年陈瘢）。形态一致：自我施加的深层抓挠——患者用自己的指甲在皮肤上反复抠出沟槽，直至皮下组织。

\noindent \textit{鉴别:} 瘢痕边缘不规则、宽度不一，为抓挠（指甲）特征，而非切割（刀片）。无提示静脉切开自杀尝试的瘢痕。形态为\textbf{强迫性非自杀性自残}，可能关联精神运动性激越发作。

\noindent \textbf{补充观察:} 最陈旧的瘢痕呈现泥土浸染性纹身重叠（愈合期间泥土颗粒混入真皮），提示伤口曾被打开并暴露于污染土壤。\textbf{患者病历中没有破伤风预防记录。}

\subsubsection*{1.7 牙齿与口腔}

\noindent \textbf{缺失12颗牙:} 上下切牙（8颗）、右下第一磨牙、左上第二前磨牙、左下尖牙、右下第三磨牙。余留牙齿呈现：

\noindent - 咬合面严重磨损（磨牙症摩擦，常见于无口腔保护接受ECT的患者）
- 4颗牙广泛龋坏并牙髓暴露
- 2颗牙未经治疗的牙折（符合ECT期间下颌收缩）

\noindent \textbf{舌:} 两侧纵向边缘各有两条约3厘米的大型线性瘢痕。为抽搐发作期间反复咬伤形成的瘢痕。电休克期间缺少口腔保护，导致患者多次咬穿自己的舌肌。瘢痕深且呈瘢痕疙瘩样——愈合组织在舌缘形成可触及的结节。

\subsubsection*{1.8 创伤性截指}

\noindent \textbf{左手小指末节指骨部分截断。} 陈旧性愈合，残端圆钝，无甲节。病历文件中没有任何可解释该截断的创伤记录。

\noindent \textit{页边注:} “二期愈合，无缝合。可能为咬伤或创伤性撕脱。患者还用牙齿拔自己的指甲（1937年4月记录）。该指节可能是在自残发作中被患者自己咬掉的。”

\subsubsection*{1.9 热烧伤}

\noindent 胸腹部前侧可见散布的\textbf{色素减退斑}，边缘不规则，直径约2至8厘米。为浅二度烧伤瘢痕，已上皮化。

\noindent \textit{分析:} 热力形态符合接触热水或加热表面。收容院护理记录提到“治疗性沐浴”做法——患者被浸入热水（40—50°C）或冷水（10—15°C）中长时间浸泡，作为“镇静”或“刺激”手段。烧伤无几何形状，提示浸烫而非接触加热物体。

\subsection*{2. 内部检查}

\subsubsection*{2.1 死因}

\noindent \textbf{心肺骤停}，继发于:
- 充血性心力衰竭
- 严重慢性营养不良（心脏恶病质）
- 支气管误吸（食物吸入下呼吸道 —— 符合神经后遗症所致吞咽困难）

\subsubsection*{2.2 心血管系统}

\noindent 心脏重量减轻（220克 —— 该体型预期：280—340克）。心肌褐色萎缩（脂褐素色素沉着，典型于长期营养不良）。心肌纤维呈胞浆空泡变性与横纹消失。主动脉轻度动脉粥样硬化，与估计年龄不符（慢性营养不良延缓了动脉粥样硬化进程）。

\subsubsection*{2.3 消化系统}

\noindent \textbf{胃:} 空虚。完全无食物内容物。胃壁变薄（平均厚度2毫米），黏膜萎缩，皱襞减少。患者死亡前至少24小时未进食。

\noindent \textbf{肠道:\textit{ 慢性吸收不良征象 —— 肠绒毛扁平缩短（绒毛近全萎缩）。存在大量 }蛔虫}（人蛔虫）——高寄生虫负荷（小肠内计数17条，升结肠4条）。每条寄生虫长15—25厘米，缠绕成团，占据约40\%的肠腔。

\noindent \textit{观察:} 高寄生虫负荷加上慢性营养不良，加剧了恶病质状态。患者与自己的寄生虫争夺那一点点营养。

\noindent \textbf{肝脏:} 重量减轻（980克），中度脂肪变性，有蛋白质—热量营养不良征象。无肝硬化或肝病征象可解释外部检查观察到的单侧黄疸。

\subsubsection*{2.4 中枢神经系统}

\noindent \textbf{脑:} 重量1,120克（在年龄与体型的预期范围内）。大脑半球对称。脑膜透明，无脑膜炎或蛛网膜下腔出血征象。

\noindent \textit{专项分析 —— 额叶切除术:} \textbf{无肉眼可见的皮质瘢痕}与经眶或前额叶切除术相符。眶额与内侧前额区被仔细解剖检查。无胶质增生、空洞或实质改变提示神经纤维切断。白质刀的进入通路（经眶顶）已被确认，但路径似乎在到达脑实质之前中断，或者器械沿一条未造成可检测脑组织损伤的解剖平面行进。

\noindent \textit{法医注:} “鉴于无脑损伤，额叶切除术可能在损伤实质前被中断，或者患者只接受了骨性穿孔（作为住院医师的‘训练’），未切断额叶纤维。病历中无记录可澄清。”

\noindent \textit{专项分析 —— 电休克:} “无宏观证据表明反复电休克治疗（估计在50年住院期间达数百次）造成可检测的脑损伤。无胶质增生、叶性萎缩或显著血管改变。”

\noindent \textit{页边注，草书笔迹，显然较新:\textit{ “脑损伤的缺失不是宽慰。是所有发现中最令人不安的。它意味着若阿金的暴力不是脑损伤造成的。暴力是}他\textit{的。或者是}别的什么东西\textit{的。那些手术没有在脑中留下痕迹——却在他全身留下了痕迹。身体的痕迹讲述了一个故事。沉默的大脑是这份报告里最刺眼的部分。 —— L.M.”}}

\subsubsection*{2.5 脊柱}

\noindent \textbf{T5压缩性骨折}（如1.3节所述）。T5—T6水平的脊髓轻度前后扁平，无显著受压。该骨折属“稳定型”，未造成神经功能缺损。\textbf{另有T3和T8较小的压缩性骨折}，在既往放射检查中未被发现。

\noindent \textit{观察:} 三处骨折椎体（T3、T5、T8）构成典型的ECT骨折三联征——发生于抽搐强直期，此时脊柱屈肌与伸肌同时收缩，产生超过骨质强度的轴向压缩力。

\subsubsection*{2.6 肌肉骨骼系统 —— 其他发现}

\noindent - \textbf{肩关节病:} 双侧肱盂关节复发性脱位（肩部多次脱位复位）。与称为“锁臂”的约束手法相关——收容院护工用来制服躁动患者的技巧，即过度外展并外旋手臂直至关节弹响。
- \textbf{头皮瘢痕:} 枕部一条6厘米线性瘢痕，符合钝器伤。可能为抽搐发作期间跌倒或工作人员身体攻击所致。
- \textbf{足部畸形:} 双侧严重拇外翻，双足第二趾叠压于拇趾之上。畸形由长期穿着不合脚的鞋（统一尺码、无鞋号——“收容院拖鞋”，由患者们在医院制鞋车间自制）造成。

\subsection*{3. 尸体上确认的手术程序汇总}

\begin{moltablethree}
\hline
\textbf{程序} & \textbf{物理证据} & \textbf{期间} \\
\hline
无麻醉电休克治疗 & 颞部烧伤；T3/T5/T8椎体骨折；下颌骨骨折；股骨骨折；舌瘢痕；牙齿缺失 & 1925—1979（54年） \\
经眶额叶切除术（弗里曼术式） & 双侧筛板孔道；结膜瘢痕 & 1950年代（无记录） \\
约束衣／长期约束 & 腕踝色素沉着带；手指挛缩；压迫性脱发 & 1920—1979（59年） \\
治疗性沐浴（浸烫） & 胸腹热烧伤瘢痕 & 1930—1960 \\
胰岛素疗法（胰岛素昏迷） & \textit{（无标志物 —— 不留物理瘢痕）} & 1940—1950 \\
湿床单约束 & 大腿与髋部真皮纤维化 & 1930—1979 \\
\hline
\end{moltablethree}

\noindent \textit{最终注（手写，事后添加）:} 尸体呈现59年的机构性干预标志：无麻醉ECT（双侧颞部烧伤，T3/T5/T8椎体骨折），经眶额叶切除术（双侧筛板穿孔），机械约束（手指挛缩，大腿真皮纤维化），治疗性沐浴（胸腹热瘢痕）。\textbf{死因:} 心肺骤停。\textbf{促成原因:} 慢性营养不良恶病质（估计体重：38公斤）。\textbf{合并症:} 高寄生虫负荷（人蛔虫，21条）。未发现任何器质性疾病标志物可解释62年的机构化收治。患者最严重的损伤不是解剖学上的。

\vspace{0.5em}\noindent\rule{\textwidth}{0.5pt}\vspace{0.5em}

\noindent 阿尔瓦罗·托雷斯医生

\noindent 法医 —— 巴尔巴塞纳法医研究所

\noindent CRM-MG 3,117

\vspace{0.5em}\noindent\rule{\textwidth}{0.5pt}\vspace{0.5em}

\noindent \textit{附言（手写，显然事后添加）:} “本报告写于1979年6月14日，掘尸之前。下一节所述的发现——掘尸与特殊发现——是在开坟并取回日记之后补充的。签字人声明，对第4节（掘尸）所述现象没有医学解释，仅限于转述掘墓人和所取回材料提供的内容。日记由法医研究所保管，附于本案卷。—— A.T.”

\noindent \textit{补充注（L.M.笔迹，2026年案卷）:} “从坟中与死者材料同一土层中取回了鸮形目（穴居小猫头鹰）骨骼碎片。若阿金墓葬中猫头鹰残骸的存在，被掘墓人记录为‘不肯离开的那只鸟’。未找到该沉积的埋藏学解释。猫头鹰骨在收容院人类墓葬中并非常见发现。”

\noindent\textit{\textbf{↪ 路线:} \rota{MEM-04} \rotasep \rota{DOC-02} \rotasep \rota{LUC-03} \rotasep \rota{CONT-03}}
